\section{A second proof of the editing lower bound}\label{sec:probmethod}
The exact disjointness of protected sets under the XOR selection is a
convenient deterministic proof. A probabilistic argument reaches the
same editing threshold using only one-coordinate uniformity. This
formulation shows which feature is intrinsic to the tree host and
which feature is special to the XOR coupling.
The marginal-exposure and union-bound method adapts the edit-distance
argument of \cite{OAI} to the smaller marker frame used here. It is
included to explain the construction and to compare it with the
deterministic packing, rather than as a new probabilistic method.

Choose a leaf label $Z$ uniformly from $\{0,\dots,m-1\}$. At every
depth $d$ select the node $\lfloor Z/2^{h-d}\rfloor$ and, in each
signed block at that node, sample one physical row and one physical
column uniformly from its $2^{h-d}$ offsets. Sample the row and
column representatives in the marker and dummy blocks uniformly
and independently of each other and of the tree representatives.
The two signed paths meet at the same selected physical leaf
because a leaf block consists of a single position.

\begin{lemma}[Single-position marginal]\label{lem:marginal}
For any fixed variable row position at depth $d$ and sign, its
probability of selection is $m^{-1}$. The same is true for a fixed
variable column position, marker position, or dummy position.
\end{lemma}
\begin{proof}
A depth-$d$ node is reached with probability $2^{-d}$.
Conditional on reaching it, a given offset is sampled with
probability $2^{-(h-d)}$. Multiplying gives $2^{-h}=m^{-1}$.
At the leaf depth the selection probability is $2^{-h}$ because
there is one representative per node. Marker and dummy blocks have
exactly $m$ positions and are sampled uniformly.
\end{proof}

Protect all selected cells for which at least one endpoint is a
marker or dummy position, together with the two selected
root-variable diagonal cells. No other variable--variable cells
are protected. In particular, the shared terminal cell, whose
value will be forced by the propagation, need not be protected.

\begin{lemma}[Two-coordinate exposure]\label{lem:twocoord}
For every fixed physical matrix cell $e$, the probability that $e$
is protected by the random selection is at most $m^{-2}$.
\end{lemma}
\begin{proof}
If either endpoint of $e$ is a marker or dummy position, that
endpoint is chosen independently of the tree-path and variable
offset choices on the other axis. Each endpoint has marginal
probability $m^{-1}$ by Lemma~\ref{lem:marginal}, and the two
events are independent, giving $m^{-2}$. If both endpoints are
variable positions at the root, their offsets are sampled
independently from the two root blocks of size $m$, giving the
same product. All other cells have protection probability zero.
\end{proof}

\begin{proposition}[Probabilistic lower bound]
\label{prop:problower}
Every edit set that turns $A_h$ into an $H_4$-free binary matrix
contains at least $m^2$ cells.
\end{proposition}
\begin{proof}
Let $E$ be an arbitrary set of changed cells with $|E|<m^2$.
By Lemma~\ref{lem:twocoord} and the union bound,
\[
 \mathbb P\bigl(E\cap\mathrm{Protected}\ne\varnothing\bigr)
 \le \sum_{e\in E}\mathbb P(e\in\mathrm{Protected})
 \le |E|m^{-2}<1.
\]
A selection therefore exists whose protected cells avoid every
edit. The forcing frames along its common path are valid and have
all twelve marker incidences unchanged. The same four implications
as in Section~\ref{sec:distance} propagate the root value one
through the plus chain and the root value zero through the minus
chain. Both implications concern the same selected terminal cell;
they are inconsistent in an $H_4$-free edited matrix. Thus $E$
cannot eliminate all copies.
\end{proof}

The probabilistic argument uses independence \emph{between axes}
whenever one position is in a marker or dummy block. It does not
require independence of the variable representatives on different
levels of the same axis, nor of the two terminal path labels.
The deterministic XOR proof goes further: each individual protected
cell identifies at most one of the $m^2$ selections. This is an exact
combinatorial packing statement, while the probabilistic proof
controls the probability of hitting each required cell. Both yield
the sharp integer threshold $m^2$.

\subsection{Boolean verification of the four implications}
The lower-right body $P$ is
$\left(\begin{smallmatrix}1&0\\1&1\end{smallmatrix}\right)$.
Each of the four forcing frames has three of its four body entries
fixed by the choice of mode or by the current inductive hypothesis.
The remaining entry is constrained by the absence of a completed
copy. The four exact incompatibilities are
\begin{align*}
 V_i^+:&\quad (x_{i-1},\alpha_i)=(1,0),\\
 V_i^-:&\quad (y_{i-1},\beta_i)=(0,1),\\
 W_{i,p_i\bmod2}^+:&\quad (\alpha_i,x_i)=(1,0),\\
 W_{i,p_i\bmod2}^-:&\quad (\beta_i,y_i)=(0,1).
\end{align*}
Each forbidden pair would complete the body $P$ and therefore
produce $H_4$ with the protected frame. The complement of each
pair yields exactly the implication in the forcing table. This
explicit truth-table reduction is independent of the arithmetic
classification certificates used for the upper bound.
